%%%%%%%%%%%%%%%%%%%%%%%%%%%%%%%%%%%%%%%%%
% Medium Length Professional CV
% LaTeX Template
% Version 2.0 (8/5/13)
%
% This template has been downloaded from:
% http://www.LaTeXTemplates.com
%
% Original author:
% Trey Hunner (http://www.treyhunner.com/)
%
% Important note:
% This template requires the resume.cls file to be in the same directory as the
% .tex file. The resume.cls file provides the resume style used for structuring the
% document.
%
%%%%%%%%%%%%%%%%%%%%%%%%%%%%%%%%%%%%%%%%%

%----------------------------------------------------------------------------------------
%   PACKAGES AND OTHER DOCUMENT CONFIGURATIONS
%----------------------------------------------------------------------------------------

\documentclass{resume} % Use the custom resume.cls style
\usepackage[left=0.7in,top=0.4in,right=0.7in,bottom=0.5in]
{geometry} % Document margins
\usepackage{url}
\usepackage{bookman}
\usepackage[T1]{fontenc}
\usepackage{hyperref}

%----------------------------------------------------------------------------------------
%   DOCUMENT START
%----------------------------------------------------------------------------------------
\name{Abhijeet Lokhande} % Your name
\designation{AI Alchemist}

\address{07480801382 \\ London, United Kingdom}
\address{\href{mailto:abhijeet_lokhande@proton.me}{abhijeet\_lokhande@proton.me} \\ https://www.linkedin.com/in/ablds/}

\begin{document}
\title{Anon Resume}
%----------------------------------------------------------------------------------------
%   EDUCATION SECTION
%----------------------------------------------------------------------------------------
\begin{rSection}{Summary}
AI Alchemist committed to continuous learning and building impactful AI solutions from fine-tuning LLMs (50K+ HuggingFace downloads) to developing AI agents and data pipelines.
Passionate about solving real-world problems through innovative machine learning approaches.
\end{rSection}

\begin{rSection}{Technical Skills}

\begin{tabular}{ @{} >{\bfseries}l @{\hspace{6ex}} l }
Proficient In   & {Machine/Deep Learning, NLP, LLMs, AI Agents, DSPy}\\
                & {Python, SQL, Django}\\
                & {Geopandas, Scikit-Learn, TensorFlow, Keras, PyTorch, XGBoost}\\
                & {Langchain, LlamaIndex, LangGraph}\\
                & {FastAPI, Test Driven Development}\\
                & {Retrieval Augmented Generation(RAG), Data Analysis}\\

                & {ETL processes, Data pipelines, MySQL, PostgreSQL, MongoDB}\\
                & {AWS, Git, GitHub, Docker} \\
                & {Numerical Optimization, Statistics, Forecasting} \\
\end{tabular}


\end{rSection}

%----------------------------------------------------------------------------------------
%   Personal Details
%----------------------------------------------------------------------------------------


\begin{rSection}{Profile Links}
\begin{tabular}{ @{} >{\bfseries}l @{\hspace{6ex}} l }
Linkedin & https://www.linkedin.com/in/ablds/  \\
Github & https://github.com/abhijeetscode \\
Hugging Face & https://huggingface.co/ab-ai
\end{tabular}

\end{rSection}


%----------------------------------------------------------------------------------------
%   WORK EXPERIENCE SECTION
%----------------------------------------------------------------------------------------
\begin{rSection}{Open Source Contributions}
\begin{itemize}
    \item Fine-tuned Large Language Models (GPT with LoRA and BERT) to detect Personally Identifiable Information (PII), with published models collectively achieving 30,000+ downloads on HuggingFace.
    \url{https://huggingface.co/ab-ai/PII-Model-Phi3-Mini} \\
    \url{https://huggingface.co/ab-ai/pii_model}
    \item Built an autonomous ReAct-based coding agent using LangGraph and the Anthropic API that generates installable, tested software packages from plain-language markdown specs, achieving 9/9 pass rate across a polyglot benchmark (Python, Go, Rust, Java, JS). \url{https://github.com/abhijeetscode/Spec-to-Code-Agent}
    \item Developed CLIP-style architecture for searching images based on
    natural language prompts using dual-encoder embeddings. \\
    \url{https://github.com/abhijeetscode/text-to-image-search}
    \item  Built a Digital Detox mobile app that helps users reduce social media usage through intentional friction and app-blocking mechanisms. \\
    \url{https://abhijeets-code.itch.io/digitaldetox}

\end{itemize}
\end{rSection}

\begin{rSection}{Work Experience}
\begin{rSubsection}{Built AI (FinTech \& AI Company)}{April-2022 - Present}{AI Alchemist (Early Stage Hire)}{London, UK}
\item Developed and deployed an autonomous AI agent that generates Discounted Cash Flow (DCF) financial models from raw financial data, reducing analyst modeling time by ~70\% and significantly improving workflow efficiency.
\item Built a production-ready LLM agent to understand and explain financial cashflows to users, with robust validation to ensure accurate and reliable outputs.
\item Developing an AI agent that interprets investment-related queries, retrieves relevant outputs from a financial modeling engine, and generates user-friendly responses with actionable insights.
\item Built a Conversational AI bot for analyzing financial documents using LLMs, RAG, embeddings, and vector database — enabled financial analysts to query multiple PDFs with natural language, cutting manual review time.
\item Developed and deployed a novel machine learning model to infer missing values in large-scale property datasets. Enabled accurate estimation of rent per square foot across London postcodes.
\item Built a data processing and ingestion pipeline by linking 80\% of property records across diverse datasets with geolocation enrichment, enabling a machine learning model to infer missing rental values by area.
\item Improved performance of the financial modelling engine by 85\% by rearchitecting the existing codebase into a graph-based system leveraging Apache Hamilton and Vectorisation techniques.
\item Delivered a 30\% productivity gain by designing and implementing an automated financial document processing pipeline leveraging AWS Textract.
\end{rSubsection}


\begin{rSubsection}{King's College London}{March 2021 - January 2022}{Research Associate}{London, UK}
\item Developed a novel Deep Learning model utilizing Graph Convolution Network and Generative Adversarial Network (GAN) techniques for Drug Discovery, resulting in a 63\% enhancement in the novelty score of
drug structures.
\item Designed and implemented a robust data preprocessing pipeline for clinical data, ensuring accuracy and completeness of the a million-row dataset.
\item Utilized Scrapy and Beautiful Soup to gather data from various websites for research purposes

\end{rSubsection}

\end{rSection}


%----------------------------------------------------------------------------------------
%   Personal Project
%----------------------------------------------------------------------------------------


\iffalse
\begin{rSection}{Relevant Project}
\begin{description}
  \item[$\cdot$] Solved infinite monkey problem using evolutionary algorithms.
  \item[$\cdot$] Automated a Pacman on a small grid using reinforcement learning.
  \item[$\cdot$] I built VAE to generate new images from the fashion MNIST dataset.
  \item[$\cdot$] I built a movie recommender system using cosine similarity.
\end{description}
\end{rSection}
\fi


%----------------------------------------------------------------------------------------
%   Achievements
%----------------------------------------------------------------------------------------
\iffalse
\begin{rSection}{Achievements}
\begin{description}
  \item[$\cdot$] Completed all level of Google FooBar.
  \item[$\cdot$] Got a certificate of appreciation from Merck \& Co.
  \item[$\cdot$] Arctic Code Vault Contributor of Github.
\end{description}
\end{rSection}
\fi
%----------------------------------------------------------------------------------------
%   Technical Experience
%----------------------------------------------------------------------------------------



%----------------------------------------------------------------------------------------
%   TECHNICAL STRENGTHS SECTION
%----------------------------------------------------------------------------------------
\begin{rSection}{Education}
% September 2020 - January 2022
\begin{rSubsection}{King's College London}{}{MSc in Data Science }{London, UK}
%\vspace{-.1cm}
%Overall GPA: 78\%
\vspace{-.65cm}
\item[]
\end{rSubsection}
\begin{rSubsection}{C-DAC: Centre for Development of Advanced Computing, India}{}{Post Graduate Diploma in Advanced Computing (PG-DAC)}{Pune, India}
%\vspace{-.1cm}
%Overall GPA: 78\%
\vspace{-.65cm}
\item[]
\end{rSubsection}
% June 2013 - June 2017
\begin{rSubsection}{University of Pune}{}{Bachelor of Engineering in Computer Science}{Pune, India}
%\vspace{-.1cm}
%Overall GPA: 78\%
\vspace{-.65cm}
\item[]
\end{rSubsection}
\end{rSection}

%----------------------------------------------------------------------------------------

\end{document}
